\section{Research Agenda}
\label{sec:research-agenda}

\subsection{Security requirements for HPC-agent platforms}
\label{sec:requirements}

The gap analysis in Section~\ref{sec:defense-gap} motivates a concrete list of properties an HPC-agent platform would need in order to make hijacked-agent incidents detectable or preventable, rather than indistinguishable from legitimate broad-account activity:

\begin{itemize}[leftmargin=*]
  \item \textbf{Task-scoped authority}, enforced independently of the account's standing access --- the platform must be able to represent and check ``permitted for this task'' as distinct from ``permitted for this account.''
  \item \textbf{Provenance-aware context handling} --- every piece of content entering an agent's context (user instruction, file content, tool output, peer-agent message) should carry a trust/source label the agent's decision process can condition on.
  \item \textbf{Egress control} for agent-to-model traffic, since data can leave through a channel --- an LLM API call --- that is authorized and encrypted by default and therefore invisible to conventional network monitoring aimed at exfiltration.
  \item \textbf{Shared-filesystem hygiene}, given that scratch and collaboration storage persists adversarial content across jobs, users, and agents in a way session-scoped web contexts do not.
  \item \textbf{Tool and module trust}, extending existing software-supply-chain practice to cover tool \emph{descriptions} and MCP manifests, not just binaries.
  \item \textbf{Scheduler-aware containment}, so that a hijacked agent cannot translate a context-level compromise into unbounded resource consumption against the user's allocation.
  \item \textbf{Cross-agent correlation}, so that staged, multi-hop exfiltration across cooperating agents is visible as a single event rather than several individually unremarkable ones.
\end{itemize}

\subsection{Toward a benchmark: HPC-AgentSec}
\label{sec:benchmark}

A threat model and a taxonomy motivate but do not by themselves measure anything. The natural next step --- already underway as a companion effort --- is an empirical benchmark that operationalizes Sections~\ref{sec:threat-model} and~\ref{sec:taxonomy} into realistic tasks, task-scoped policies, injected attack cases, deterministic security oracles, and joint utility/security metrics, following the general methodology established by AgentDojo \cite{Debenedetti2024AgentDojo} but built around HPC-native semantics: Slurm-mediated scheduler actions, project-scoped filesystem policy, and scientific-integrity oracles that check for silent parameter, filter, or provenance tampering rather than only data exposure. That effort (working name \textbf{HPC-AgentSec}; see the companion benchmark paper draft and development plan in the same project repository) is scoped to answer, empirically, the questions this position paper raises analytically: how often current agents follow adversarial HPC context, which infrastructure surfaces are highest-risk, and which existing or proposed controls actually reduce risk without destroying utility.

\subsection{Open questions}
\label{sec:open-questions}

Beyond the benchmark itself, several questions remain open and are, in our view, worth independent study:

\begin{itemize}[leftmargin=*]
  \item How should ``task scope'' be authored and validated when it is inherently somewhat subjective, and how much does that subjectivity affect the reliability of security measurements built on it?
  \item Do defenses that work well in web/enterprise agent settings (structured provenance labels, tool allowlists) transfer to HPC's scheduler and filesystem semantics, or do they need HPC-specific redesign?
  \item What is the right locus of enforcement --- the agent framework, the scheduler, the filesystem, or a dedicated runtime monitor --- for task-scoped policy, given that HPC sites differ widely in what they can modify?
  \item As multi-agent scientific workflows become more common, how should cross-agent trust be represented so that staged exfiltration is detectable without requiring every agent to fully trust every other agent's output?
\end{itemize}
